\subsection{媒介変数で表された曲線}

円や，縦向きの接線を持つ曲線は，一つの関数 $y=f(x)$ で扱いにくい場合がある。そこで
\[
x=x(t),\qquad y=y(t),\qquad \alpha\leq t\leq\beta
\]
と表す。$t$ を少し増やしたとき，横の変化はほぼ $x'(t)\Delta t$，縦の変化はほぼ $y'(t)\Delta t$ だから，線分の長さはほぼ
\[
\sqrt{x'(t)^2+y'(t)^2}\,\Delta t
\]
である。$x,y$ が連続微分可能なら，導関数の連続性によりこの近似の誤差は区間全体で制御でき，折れ線の極限として
\begin{equation}\label{式:媒介曲線の弧長}
L=\int_\alpha^\beta\sqrt{x'(t)^2+y'(t)^2}\dd t
\end{equation}
を得る。ここでは近似誤差の評価の詳細は省き，この弧長公式を用いる。途中で速度に当たる平方根が零になる点があっても，この式自体は使える。

ただし，この積分は $t$ に沿って進んだ長さを数える。同じ線分や曲線を二度たどれば二度数える。たとえば $x=\cos t$，$y=\sin t$ を $0\leq t\leq4\pi$ で積分すると $4\pi$ だが，点の集合としての単位円の周長は $2\pi$ である。図形そのものの長さを求めるなら，一度だけたどる区間を選ぶ。

媒介変数は面積にも使える。$x(t)$ が単調増加し，$y(t)\geq0$ で，曲線が $x$ の関数として表される場合は，置換 $\dd x=x'(t)\dd t$ により
\begin{equation}\label{式:媒介曲線の面積}
S=\int_\alpha^\beta y(t)x'(t)\dd t
\end{equation}
となる。$x(t)$ が途中で折り返すなら，区間を分けて曲線が囲む部分を確認する。この式に絶対値を付けるだけで常に図形の面積になるわけではない。帯が重なれば，同じ面積を二度数えるからである。

\begin{example}{一つの媒介表示から面積と長さを求める}
$a>0$ とする。曲線
\[
x=a(t-\sin t),\qquad y=a(1-\cos t),\qquad 0\leq t\leq2\pi
\]
の長さと，この曲線と $x$ 軸に囲まれた部分の面積を求めよ。
\end{example}

$x'(t)=a(1-\cos t)$，$y'(t)=a\sin t$ である。まず弧長の平方根を計算すると
\[
\begin{aligned}
x'(t)^2+y'(t)^2
 &=a^2\{(1-\cos t)^2+\sin^2t\}\\
 &=2a^2(1-\cos t)=4a^2\sin^2\frac t2.
\end{aligned}
\]
平方根は $2a|\sin\frac t2|$ である。この区間では $0\leq\frac t2\leq\pi$ なので正弦は非負であり，
\[
L=2a\int_0^{2\pi}\sin\frac t2\dd t
 =\left[-4a\cos\frac t2\right]_0^{2\pi}=8a.
\]
平方根を取る際に絶対値が必要かどうかを，変数の範囲で確認している。

一方，$x'$ は端点以外で正なので，$x$ は区間全体で単調増加する。また $y\geq0$ で，両端では $y=0$ である。よって面積は
\[
\begin{aligned}
S&=a^2\int_0^{2\pi}(1-\cos t)^2\dd t\\
 &=a^2\int_0^{2\pi}\left(\frac32-2\cos t+\frac12\cos2t\right)\dd t
 =3\pi a^2.
\end{aligned}
\]
弧長では二方向の変化から距離を作り，面積では高さに横の幅を掛けた。同じ媒介表示でも，求める量に応じて積分の中身が異なる。

